% ==========================================================================
% Präambel: Pakete, Schriftarten, Layout und Theoreme
% Angelehnt an das Design der Vorlesungsmitschrift Lineare Algebra 1
% ==========================================================================

\usepackage[a4paper, margin=2.2cm]{geometry}
\usepackage{fontspec}

% Sprach- und Schrifteinstellungen (geeignet für LuaLaTeX & XeLaTeX)
\usepackage[english, ngerman]{babel}
\babelfont{rm}{Noto Sans}

% Mathematische Pakete
\usepackage{amsmath, amssymb, amsthm}
\usepackage{booktabs}
\usepackage{enumitem}
\usepackage{xcolor}
\usepackage{hyperref}

% Farbpalette
\definecolor{primary}{RGB}{20, 50, 110}
\definecolor{secondary}{RGB}{90, 105, 125}
\definecolor{boxbg}{RGB}{246, 248, 252}
\definecolor{bordercol}{RGB}{210, 220, 235}

% Theorem-Definitionen
\theoremstyle{definition}
\newtheorem{definition}{Definition}[section]
\newtheorem{beispiel}[definition]{Beispiel}
\newtheorem{bemerkung}[definition]{Bemerkung}

\theoremstyle{plain}
\newtheorem{satz}[definition]{Satz}
\newtheorem{lemma}[definition]{Lemma}

% Listen-Formatierung
\setlist[itemize]{label=--, leftmargin=1.8em, itemsep=2pt}
\setlist[enumerate]{leftmargin=1.8em, itemsep=2pt}

% Hyperlink-Stile
\hypersetup{
    colorlinks=true,
    linkcolor=primary,
    urlcolor=primary,
    pdftitle={Grundlagen der Mathematik I - Vorlesungsmitschrift},
    pdfauthor={Soroush}
}

% Standard-Mengensymbole als handliche Makros
\newcommand{\N}{\mathbb{N}}
\newcommand{\Z}{\mathbb{Z}}
\newcommand{\Q}{\mathbb{Q}}
\newcommand{\R}{\mathbb{R}}
\newcommand{\C}{\mathbb{C}}
\newcommand{\K}{\mathbb{K}}
\newcommand{\F}{\mathbb{F}}
\newcommand{\Pot}{\mathcal{P}}
\newcommand{\xor}{\mathbin{\dot{\lor}}}

% Vorlesungs-Marker: Datum und zugehörige handschriftliche Quelle
\newcommand{\vorlesungsdatum}[2]{%
    \vspace{0.6em}%
    \noindent\textcolor{secondary}{\small\textbf{Vorlesung vom #1} \hfill \textit{Handschrift: \texttt{#2}}}\par%
    \vspace{0.4em}%
}
